\documentclass[conference]{IEEEtran}
\IEEEoverridecommandlockouts
% The preceding line is only needed to identify funding in the first footnote. If that is unneeded, please comment it out.
% \usepackage{cite}
\usepackage{amsmath,amssymb,amsfonts}
\usepackage{algorithmic}
\usepackage{graphicx}
% \usepackage{textcomp}
\usepackage{xcolor}
\usepackage{booktabs}
\usepackage{multirow} 
\usepackage{url}
\usepackage{listings}
\usepackage{array} % Required for column width adjustments
\usepackage{tabularx}
\usepackage{balance}
\usepackage{colortbl} 
% \usepackage{footnote}
\usepackage{natbib} 
\usepackage{adjustbox}
\usepackage{makecell}
\usepackage{hyperref}

\def\BibTeX{{\rm B\kern-.05em{\sc i\kern-.025em b}\kern-.08em
    T\kern-.1667em\lower.7ex\hbox{E}\kern-.125emX}}
\begin{document}

\title{
LLMs at the Edge: Performance and Efficiency Evaluation with Ollama on Diverse Hardware
}

\author{Anonymous Submission}

\iffalse
\author{\IEEEauthorblockN{1\textsuperscript{st} Given Name Surname}
\IEEEauthorblockA{\textit{dept. name of organization (of Aff.)} \\
\textit{name of organization (of Aff.)}\\
City, Country \\
email address or ORCID}
\and
\IEEEauthorblockN{2\textsuperscript{nd} Given Name Surname}
\IEEEauthorblockA{\textit{dept. name of organization (of Aff.)} \\
\textit{name of organization (of Aff.)}\\
City, Country \\
email address or ORCID}
\and
\IEEEauthorblockN{3\textsuperscript{rd} Given Name Surname}
\IEEEauthorblockA{\textit{dept. name of organization (of Aff.)} \\
\textit{name of organization (of Aff.)}\\
City, Country \\
email address or ORCID}
\and
\IEEEauthorblockN{4\textsuperscript{th} Given Name Surname}
\IEEEauthorblockA{\textit{dept. name of organization (of Aff.)} \\
\textit{name of organization (of Aff.)}\\
City, Country \\
email address or ORCID}
\and
\IEEEauthorblockN{5\textsuperscript{th} Given Name Surname}
\IEEEauthorblockA{\textit{dept. name of organization (of Aff.)} \\
\textit{name of organization (of Aff.)}\\
City, Country \\
email address or ORCID}
\and
\IEEEauthorblockN{6\textsuperscript{th} Given Name Surname}
\IEEEauthorblockA{\textit{dept. name of organization (of Aff.)} \\
\textit{name of organization (of Aff.)}\\
City, Country \\
email address or ORCID}
}
\fi

\maketitle

\begin{abstract}
%background
Large Language Models (LLMs) have revolutionized natural language processing, enabling a wide range of applications. However, deploying LLMs on edge devices remains challenging due to computational and memory constraints. 
%objectives
This study addresses these challenges by leveraging Ollama, an open-source framework optimized for efficient local LLM deployment, offering insights into optimizing edge AI.
%method (how and what we do)
We assess the performance and efficiency of open-source LLMs from the Qwen2.5 and Llama3 families, spanning 0.5B to 90B parameters, leveraging the Ollama framework. Our evaluation encompasses diverse hardware platforms, including Apple Silicon (M3/M4 Max), NVIDIA GPUs (RTX A6000, RTX 4090), and edge devices such as the NVIDIA Jetson AGX Orin, using both 4-bit quantization and FP16 precision.
%results
Key findings include: 
%(1) Ollama’s \texttt{Q4\_K\_M} quantization improves F1 scores by 3.28\% for 8B models and 0.32\% for 70B models compared to 4-bit bitsandbytes quantization;
% (1) Ollama’s \texttt{Q4\_K\_M} quantization improves F1 scores for smaller models compared to 4-bit bitsandbytes quantization, with a noticeable increase for models of moderate size and a smaller improvement for very large models;
(1) Ollama’s \texttt{Q4\_K\_M} quantization outperforms the 4-bit bitsandbytes baseline, matches or exceeds FP16 performance, and achieves faster speeds with lower memory usage.
(2) Windows Subsystem for Linux (WSL) achieves significantly better throughput and lower evaluation times compared to native Windows, with up to 21× speed improvements for small models; 
%(3) medium-sized models (7B–32B) strike an optimal balance between performance and efficiency, exemplified by Qwen2.5:32B’s 96.75\% F1 score and practical inference speeds;
%(4) larger models exhibit greater cross-platform consistency, with F1 score variance decreasing from 0.26--0.9\% in small models to 0.07--0.23\% in large models.
(3) Medium-sized models (7B–32B) strike an optimal balance between performance and efficiency, as demonstrated by Qwen2.5’s high F1 score and practical inference speeds;
(4) Larger models show greater cross-platform consistency,  compared to smaller models.
%conclusion
These findings highlight the potential of robust quantization techniques and platform-specific optimizations for enabling scalable, efficient, and accessible LLM deployments in resource-constrained environments.

\end{abstract}

\begin{IEEEkeywords}
Large Language Models (LLMs), edge AI, Ollama framework, hardware benchmarking, model quantization,  performance analysis, inference throughput, cross-platform evaluation
\end{IEEEkeywords}

\section{Introduction}

The transformative impact of Large Language Models (LLMs) on natural language processing is undeniable, enabling diverse applications ranging from text summarization to complex question answering~\cite{zheng2024review}. However, deploying these sophisticated models on resource-constrained edge devices presents significant challenges due to limited computational power, memory, and energy resources~\cite{dhar2024empirical, zheng2024review}. This necessitates efficient deployment strategies that balance performance with resource constraints, particularly in edge AI environments.

The growing demand for edge-based LLMs is driven by their advantages, including reduced latency, enhanced privacy, and offline operation capabilities. These attributes make edge LLMs indispensable for critical infrastructure and real-time applications~\cite{zheng2024review}. To address these challenges, a range of optimization techniques has emerged, such as quantization~\cite{kim2024shortened, frantar2022optimal}, pruning~\cite{kim2024shortened, frantar2022optimal,yang2022distilling}, knowledge distillation~\cite{gu2023minillm, yang2022distilling}, and low-rank approximation~\cite{ji2024agile}. These techniques enable efficient LLM deployment across heterogeneous edge devices, from smartphones and wearables to IoT sensors and embedded systems~\cite{zheng2024review}.

In this context, Ollama emerges as a lightweight and extensible framework designed to simplify local LLM deployment and execution\footnote{\url{https://ollama.com/}}. By providing a straightforward API for creating and managing models, Ollama addresses the complexities often associated with LLM deployment while delivering enhanced transparency, privacy, and performance. Its cross-platform compatibility across macOS, Linux, and Windows empowers organizations to efficiently leverage existing hardware infrastructure for AI applications.

This paper evaluates open-source LLMs for edge AI applications using Ollama. Specifically, we examine models from the \textit{Qwen2.5}~\cite{yang2024qwen2} and \textit{Llama3}~\cite{grattafiori2024llama3} families, covering a range of model sizes (0.5B-90B parameters) and precisions (4-bit quantization and FP16). The evaluation encompasses a diverse set of hardware configurations, including Apple M3/M4 laptops, NVIDIA RTX A6000/4090 GPUs, and the NVIDIA Jetson AGX Orin.

Our main contributions are as follows:

\begin{itemize}
    % \item The study leverages Ollama, an open-source framework optimized for efficient local deployment of LLMs on edge devices. We analyze Ollama's performance in deploying LLMs across six hardware platforms, offering insights into performance trade-offs for edge AI applications.

    % \item %Windows Subsystem for Linux (WSL) delivers significantly higher inference speeds than native Windows, with throughput differences up to 21× for smaller models. Running Ollama in WSL is crucial for optimal efficiency.
    % The study evaluates the performance and efficiency of open-source LLMs from the Qwen2.5 and Llama3 families across a range of hardware platforms, including Apple Silicon, NVIDIA GPUs, and edge devices like the NVIDIA Jetson AGX Orin.

    \item This study leverages Ollama, an open-source framework, to evaluate the performance and efficiency of Qwen2.5 and Llama3 models across six hardware platforms, including Apple Silicon, NVIDIA GPUs, and edge devices like the NVIDIA Jetson AGX Orin, providing insights into performance trade-offs for edge AI applications.
    \item The study finds that running Ollama on Windows Subsystem for Linux (WSL) delivers significantly higher inference speeds compared to native Windows, with throughput improvements of up to 21× for smaller models, making WSL essential for achieving optimal efficiency.
    % \item %Ollama's \texttt{Q4\_K\_M} quantization outperforms bitsandbytes' 4-bit implementation, with \texttt{llama3.1:8b} achieving an F1 score of 92.78\% ($\pm$0.50\%) across six platforms, surpassing the baseline of 89.50\%. Similarly, \texttt{llama3.1:70b} achieves 95.83\% ($\pm$0.07\%), exceeding 95.51\%. 
    % This study discovers that Ollama's \texttt{Q4\_K\_M} quantization significantly improves F1 scores for smaller models compared to 4-bit bitsandbytes quantization, with notable improvements for moderate-sized models. Larger models exhibit greater stability, with F1 score variance decreasing from 0.26--0.9\% (small models) to 0.07--0.23\% (large models). Medium-sized models (7B--32B) strike an optimal balance between performance and efficiency, with \texttt{qwen2.5:32B} achieving a 96.75\% F1 score and practical inference speeds.
    \item This study reveals that Ollama's \texttt{Q4\_K\_M} quantization delivers outstanding performance and efficiency, surpassing the 4-bit implementation baseline of bitsandbytes while achieving comparable or superior performance to FP16 with significantly lower memory requirements and faster inference speeds. Medium-sized models (7B--32B) strike an optimal balance between performance and efficiency, with \texttt{qwen2.5:32B} achieving a 96.75\% F1 score and practical inference speeds.

    \item This study develops an evaluation framework measuring F1 scores, throughput, memory usage, and latency across edge platforms. Upon acceptance, we will make our code publicly available, offering a reusable resource for future research.
\end{itemize}

\iffalse
Our main contributions are as follows:

\begin{itemize}
    \item 
    We conduct a rigorous analysis of Ollama's capabilities in deploying LLMs across six diverse hardware platforms, offering detailed insights into performance trade-offs for edge AI applications.

    \item 
    Our analysis identifies size-dependent performance stability, with larger models exhibiting reduced variance in F1 scores (0.07--0.23\%) across platforms, compared to smaller models (0.26--0.9\%).

    \item  
    We discover that medium-sized models (7B--32B) strike the optimal balance between performance and efficiency for edge deployment. Notably, the Qwen2.5:32B model achieves a 96.76\% F1 score while maintaining practical inference speeds across various hardware platforms.

    \item We find that, while F1 scores are nearly identical across all models, inference speed in Windows Subsystem for Linux (WSL) surpasses that of native Windows, with differences up to 21x. Given that most consumer-grade NVIDIA GPUs operate on Windows, running Ollama in WSL is crucial for optimal efficiency.

    \item 
    We develop a comprehensive evaluation framework to measure key metrics such as F1 scores, throughput, memory usage, and latency across diverse edge platforms. This framework provides actionable insights for optimizing edge AI deployments and serves as a reusable resource for future research.
\end{itemize}
\fi

\section{Related Work}

The deployment of LLMs on edge devices has garnered increasing attention, driven by the rising demand for localized, low-latency, and privacy-preserving AI capabilities. Existing literature highlights several advancements and challenges in the efficient deployment and inference of LLMs on resource-constrained edge hardware.

Significant progress has been made in optimizing LLMs for edge deployment using techniques such as quantization, pruning, and knowledge distillation. Quantization methods, including INT4 and FP8 precision, have shown promise in reducing model size and computational demands while maintaining performance~\cite{zheng2024review, dhar2024empirical, chitty2024llm}. Pruning techniques—both structured and unstructured—streamline models by removing redundant parameters, enhancing their feasibility for edge applications~\cite{zheng2024review}. Additionally, knowledge distillation facilitates the transfer of capabilities from large models to smaller, edge-compatible versions, preserving task-specific performance~\cite{zheng2024review}. The application of these techniques to transformer-based architectures remains an active research area, as their high-dimensional operations pose unique challenges~\cite{zheng2024review, dhar2024empirical}.

% Runtime optimizations are crucial for achieving efficient LLM inference on edge devices. Framework-level solutions, such as TensorRT-LLM, vLLM, Ollama, and llama.cpp, provide low-latency and high-throughput inference by leveraging hardware accelerators and optimized memory management~\cite{chitty2024llm}. Techniques like KV caching and dynamic batching significantly enhance throughput by minimizing redundant computations and maximizing hardware utilization~\cite{chitty2024llm}. Furthermore, hardware-software co-design approaches have demonstrated potential in aligning inference frameworks with device-specific capabilities, leading to improved energy efficiency and reduced response times~\cite{zheng2024review, chitty2024llm}.

Despite advancements, several challenges persist in deploying LLMs on edge devices. Limited memory and computational resources on devices like Raspberry Pi and Jetson AGX Orin constrain the feasible model sizes and require aggressive compression strategies~\cite{dhar2024empirical, chitty2024llm}. Trade-offs between model precision and accuracy, as well as the latency and energy costs associated with different hardware platforms, necessitate careful consideration in model and platform selection~\cite{dhar2024empirical, chitty2024llm}. Additionally, ensuring compatibility across heterogeneous edge environments remains a key obstacle~\cite{zheng2024review}.

% Edge-deployed LLMs have been explored across various domains, including real-time translation, conversational AI, and on-device personal assistants. Their ability to operate without persistent cloud connectivity makes them suitable for privacy-sensitive and latency-critical applications~\cite{zheng2024review, dhar2024empirical}. Industrial use cases, such as robotics and autonomous vehicles, also benefit from edge LLMs for real-time decision-making and contextual analysis~\cite{zheng2024review, dhar2024empirical}.

% This paper extends the existing body of work by evaluating a diverse range of open-source LLMs using the Ollama framework on multiple edge platforms. By analyzing trade-offs between model size, precision, and hardware configurations, our study provides actionable insights into optimizing LLM deployments for edge AI applications.

This research presents LLMs for Edge AI, extending the existing body of work by evaluating a diverse range of open-source LLMs using the Ollama framework across multiple edge platforms.

\section{%Methodology
LLMs for Edge AI}

\subsection{Dataset and Prompt Engineering for LLMs at the Edge Evaluation}

The dataset used in this study builds upon the \textit{Global Maritime Risk Incident Dataset (GMRID)} introduced in~\cite{huang2024automating}. Specifically, we extend the \textit{GMRID v2} dataset by incorporating a new label column, \textit{gpt-4o\_label}, formatted in JSON and generated using the OpenAI's GPT-4o model, as illustrated in List~\ref{lst:exemplars-prompt}. Additionally, we enhance the dataset by merging the \textit{Headline} and \textit{Details} columns into a new composite column, \textit{Headline\_Details}, which serves as input for LLM prompts. This extended version of the dataset is referred to as the \textit{GMRID v3} dataset.

We use this dataset because it is relatively new and specifically developed for the maritime domain, an area where opportunities for exploration using LLMs are still limited, which ensures the reliability of the evaluation in this study.

Using the \textit{GMRID v3} dataset, we developed a prompting strategy building upon the few-shot techniques introduced in~\cite{huang2024automating}. 
%In addition to prompting LLMs for classification into primary categories, we extend the prompt to include the generation of specific tags, thereby providing fine-grained categorical information. 
The system prompt template is provided in Listing~\ref{lst:llm-classification-prompt}.

\vspace{2.0\baselineskip}
\begin{lstlisting}[caption={System Prompt for LLM Classification}, 
                   label={lst:llm-classification-prompt}, 
                   breaklines=true, 
                   basicstyle=\scriptsize\ttfamily, 
                   xleftmargin=0pt, 
                   framexleftmargin=0pt, 
                   framesep=0pt, 
                   aboveskip=0pt, 
                   belowskip=0pt, 
                   showspaces=false, 
                   showstringspaces=false,
                   keepspaces=true,
                   breakindent=0pt,
                   literate={“}{``}1 {”}{''}1 {‘}{`}1 {’}{'}1 {–}{-}1,
                   inputencoding=utf8]

Task: Classify Inputs into Predefined Categories
Your primary objective is to analyze the given input and assign it to one of the predefined categories: {categories_list}. Evaluate the content carefully and use the defining characteristics of each category to ensure an accurate classification.

[Additional details omitted for clarity]

Instructions for Output:

1. Once the category is identified, provide "specific tags" by selecting from the list corresponding to the identified category, as defined in the JSON.
2. Ensure the selected "specific tags" accurately reflect the details and context of the input.

Output Format:

Return your classification in the following JSON format:
{
  "category": "<Selected Category>",
  "specific_tags": ["<Selected Tag 1>", ...]
}
\end{lstlisting}


% Guidelines:
% 1. Understand the Categories:
% Each category has specific attributes that distinguish it. Familiarize yourself with these attributes by referring to the category descriptions provided in the JSON below. Use these details to guide your classification:

% {categories_json}

% 2. Contextual Analysis:
% Consider the broader context of the input. If an input could potentially fit into multiple categories, select the one that most closely aligns with its primary intent or focus.
% 3. Handling Ambiguity:
% For ambiguous inputs or those that do not clearly align with any category, choose the category that most closely matches the content provided.
% 4. Ensure Accuracy and Consistency:
% Strive for consistent and accurate classifications. Avoid arbitrary or random assignments.
% 5. Provide Feedback:
% If the input cannot be classified into any of the given categories, classify it as "Others."

This system prompt template extends the one in~\cite{huang2024automating} by:
\begin{itemize}
    \item Adding instructions to output specific tags.
    \item Mandating JSON as the output format.
\end{itemize}
Correspondingly, exemplars are retrieved from the \textit{GMRID v3} training set and formatted as shown in Listing~\ref{lst:exemplars-prompt}.

\begin{lstlisting}[caption={Exemplars Used in Few-Shot Prompting}, 
                   label={lst:exemplars-prompt}, 
                   breaklines=true, 
                   basicstyle=\scriptsize\ttfamily, 
                   xleftmargin=0pt, 
                   framexleftmargin=0pt, 
                   framesep=0pt, 
                   aboveskip=0pt, 
                   belowskip=0pt, 
                   showspaces=false, 
                   showstringspaces=false,
                   keepspaces=true,
                   breakindent=0pt]
Example Inputs and Outputs:
- Input:

[Additional details omitted for clarity]

- Output:
{
  "category": "Weather",
  "specific_tags": ["Severe Winds"]
}
\end{lstlisting}
% Local sources reported that operations at Pier 1 and 2 container terminals at the Port of Durban have been suspended due to strong winds...

To improve few-shot learning, we ensure each sample belongs to a different category until the shot number exceeds the total number of categories, which is eight. Additionally, we introduced more granular shot numbers: 0, 1, 2, 4, 8, and 10, replacing the original sequence of 0, 1, 3, 5, and 10.

\subsection{Evaluation Metrics}

To evaluate the performance of each model, we calculated predictive accuracy on the held-out test data using the weighted F1 score. The category elements retrieved from the LLM outputs and the \textit{gpt-4o\_label} column were used to calculate the F1 scores. This metric provides a balanced measure of the model's performance by considering both precision and recall, offering a comprehensive perspective on the model's ability to correctly classify data.

To assess the inference speed of different models and configurations, we used the following metrics:

\begin{enumerate}
    \item \textbf{Mean Evaluation Time}: This metric measures the average time (in seconds) required for the LLM to complete a classification task. It is abbreviated as \textbf{Time (s)} and is calculated as:
    \[
    \text{Mean Evaluation Time} = \frac{\text{Total Evaluation Time (s)}}{\text{Size of Test Dataset}}
    \]

    \item \textbf{Throughput}: This metric indicates the number of tokens processed per second. It is abbreviated as \textbf{T-put (t/s)} and reflects the processing efficiency of a model. Following the definition introduced by~\cite{chitty2024llm}, it is calculated as:
    \[
    \text{Throughput} = \frac{\text{Total Input \& Output Tokens (t)}}{\text{Total Evaluation Time (s)}}
    \]
\end{enumerate}

These metrics provide a comprehensive assessment of both predictive performance and computational efficiency. The weighted F1 score captures classification accuracy, while throughput and mean evaluation time emphasize the practical feasibility of deploying these models on various hardware configurations.

%\section{Experimental Setup}

\subsection{Model Selection and Specifications}

The evaluated LLMs are detailed in Table~\ref{tab:model_specs}. This study focuses on two leading open-source model families: Qwen2.5 and Llama3, both offering a diverse range of model sizes and configurations. For models with 14 billion parameters or fewer, we evaluate both quantized (\texttt{Q4\_K\_M}) and FP16 versions, enabling an exploration of trade-offs between precision and computational efficiency. For larger models, only quantized versions are assessed due to memory constraints. 

All model configurations adhere to Ollama's standardized naming conventions to ensure consistency and reproducibility in deployment and evaluation. This diverse set of architectures and scales facilitates a thorough analysis of model performance and feasibility in edge AI applications.

\begin{table}[htbp]
\centering
\begin{adjustbox}{width=0.3\textwidth}
\begin{tabular}{llr}
\toprule
\textbf{Name} & \textbf{Size} & \textbf{Precision} \\ 
\midrule
qwen2.5:0.5b & 397 MB & \texttt{Q4\_K\_M} \\ 
qwen2.5:0.5b-instruct-fp16 & 994 MB & FP16 \\ 
\cmidrule(lr){1-3}
llama3.2:1b & 1.3 GB & \texttt{Q4\_K\_M} \\ 
llama3.2:1b-instruct-fp16 & 2.5 GB & FP16 \\ 
\cmidrule(lr){1-3}
qwen2.5:1.5b & 986 MB & \texttt{Q4\_K\_M} \\ 
qwen2.5:1.5b-instruct-fp16 & 3.1 GB & FP16 \\ 
\cmidrule(lr){1-3}
llama3.2:3b & 2.0 GB & \texttt{Q4\_K\_M} \\ 
llama3.2:3b-instruct-fp16 & 6.4 GB & FP16 \\ 
\cmidrule(lr){1-3}
qwen2.5:3b & 1.9 GB & \texttt{Q4\_K\_M} \\ 
qwen2.5:3b-instruct-fp16 & 6.2 GB & FP16 \\ 
\cmidrule(lr){1-3}
qwen2.5:7b & 4.7 GB & \texttt{Q4\_K\_M} \\ 
qwen2.5:7b-instruct-fp16 & 15 GB & FP16 \\ 
\cmidrule(lr){1-3}
llama3.1:8b & 4.9 GB & \texttt{Q4\_K\_M} \\ 
llama3.1:8b-instruct-fp16 & 16 GB & FP16 \\ 
\cmidrule(lr){1-3}
llama3.2-vision:11b & 7.9 GB & \texttt{Q4\_K\_M} \\ 
llama3.2-vision:11b-instruct-fp16 & 21 GB & FP16 \\ 
\cmidrule(lr){1-3}
qwen2.5:14b & 9.0 GB & \texttt{Q4\_K\_M} \\ 
qwen2.5:14b-instruct-fp16 & 29 GB & FP16 \\ 
\cmidrule(lr){1-3}
qwen2.5:32b & 19 GB & \texttt{Q4\_K\_M} \\ 
\cmidrule(lr){1-3}
llama3.1:70b & 42 GB & \texttt{Q4\_K\_M} \\ 
\cmidrule(lr){1-3}
llama3.3:70b & 42 GB & \texttt{Q4\_K\_M} \\ 
\cmidrule(lr){1-3}
qwen2.5:72b & 47 GB & \texttt{Q4\_K\_M} \\ 
\cmidrule(lr){1-3}
llama3.2-vision:90b & 54 GB & \texttt{Q4\_K\_M} \\ 
\bottomrule
\end{tabular}
\end{adjustbox}
\caption{Specifications of the evaluated Qwen2.5 and Llama3 Ollama models. The size values represent the file sizes of the model weights. The precision values \texttt{Q4\_K\_M} and FP16 correspond to quantized 4-bit precision and half-precision floating-point formats, respectively.}
\label{tab:model_specs}
\end{table}

\subsection{Hardware Infrastructure and Test Environments}

To evaluate the selected Ollama models, we conducted performance testing across six hardware systems and eight configurations, covering diverse computing environments:

\begin{enumerate}
    \item \textbf{Apple M4 Max MacBook Pro:} Running macOS Sequoia 15.1, this system features the M4 Max chip with 48\,GB of unified memory, optimized for high-performance machine learning tasks. Throughout this paper, this configuration is referred to as \textbf{M4 Max}.
    
    \item \textbf{Apple M3 Max MacBook Pro:} Operating on macOS Sequoia 15.1, this system uses the M3 Max chip with 96\,GB of unified memory, designed for energy-efficient, high-performance workloads. Referred to as \textbf{M3 Max}.
    
    \item \textbf{NVIDIA RTX A6000 Ada Generation GPU:} Equipped with 48\,GB of GDDR6 ECC memory, 18,176 CUDA cores, and 568 Tensor cores. Tested on Windows 11 Pro (\textbf{RTX A6000 [Windows]}) and within the Windows Subsystem for Linux (WSL) (\textbf{RTX A6000}).
    
    \item \textbf{NVIDIA GeForce RTX 4090 GPU:} Running on Ubuntu 22.04 via WSL, this GPU offers 24\,GB of GDDR6X memory, 16,384 CUDA cores, and 512 Tensor cores. Referred to as \textbf{RTX 4090}.
    
    \item \textbf{NVIDIA GeForce RTX 4090 Laptop GPU:} This high-performance mobile GPU features 16\,GB of GDDR6 memory, 9,728 CUDA cores, and 304 Tensor cores. Tested on Windows 11 Pro (\textbf{RTX 4090 Laptop [Windows]}) and WSL (\textbf{RTX 4090 Laptop}).
    
    \item \textbf{NVIDIA Jetson AGX Orin Developer Kit:} Running Ubuntu 20.04, this system combines a 12-core Arm Cortex-A78AE CPU with an Ampere GPU (2,048 CUDA cores and 64 Tensor cores) and 64\,GB of memory, delivering up to 275\,TOPS for edge AI applications. Referred to as \textbf{Jetson AGX Orin}.

\end{enumerate}

% Our evaluation infrastructure spans from high-performance workstations to edge devices, ensuring a comprehensive assessment across both consumer-grade and specialized AI platforms.

The testing environment leverages Ollama 0.5.4 for executing LLMs on local GPUs, seamlessly integrated with a custom Python evaluation framework.
% built on LangChain\footnote{\url{https://www.langchain.com}}. 
This setup enables efficient utilization of computational resources, streamlining evaluation processes while maintaining adaptability across diverse computing environments.

\section{Experimental Results and Discussions}

\subsection{Improved Few-Shot Learning with \textit{GMRID v3}}

To evaluate the efficacy of the new prompting approach and the updated dataset \textit{GMRID v3}, we assessed the F1 scores of several language models: Llama-3.1-8B, Llama-3.1-70B, GPT-4o-mini, and GPT-4o. These models were also analyzed in ~\cite{huang2024automating}. The evaluations of the GPT-4o and GPT-4o-mini models were conducted via the OpenAI API, while the Llama-3.1 models were evaluated on Nvidia H100 NVL GPUs using the Hugging Face Transformers library~\cite{wolf2020transformers}, with 4-bit quantization implemented through the bitsandbytes library\footnote{\url{https://github.com/bitsandbytes-foundation/bitsandbytes}}. The resulting F1 scores for the models tested on \textit{GMRID v3}, alongside the existing results for \textit{GMRID v2}, are presented in Table~\ref{tab:few_shot_results}.

Our analysis reveals a clear improvement in the top F1 scores achieved by each model when using \textit{GMRID v3} compared to \textit{GMRID v2}. For instance, the peak performance of Llama-3.1-8B increased from 87.08\% on \textit{GMRID v2} to 89.50\% on \textit{GMRID v3}. Similarly, 
%GPT-4o-mini reached a higher top score of 94.92\% with \textit{GMRID v3}, up from 93.86\% with \textit{GMRID v2}. The larger 
Llama-3.1-70B model achieved a maximum F1 score of 95.51\% on \textit{GMRID v3}, surpassing its peak of 94.01\% on \textit{GMRID v2}. 
%Furthermore, GPT-4o demonstrated a substantial improvement in its best performance, reaching 99.74\% on \textit{GMRID v3} compared to 98.76\% on \textit{GMRID v2}. 
This consistent increase in the maximum achievable F1 scores across different models highlights the effectiveness of the new prompting approach and the refined dataset. These \textit{GMRID v3} results will also serve as a baseline reference for future experiments using Ollama, where additional models will be evaluated.

\begin{table*}[htbp]
\centering
\begin{adjustbox}{width=0.6\textwidth}
\begin{tabular}{llccccccccc}
\toprule
\textbf{Model} & Dataset & 0-shot & 1-shot & 2-shot & 3-shot & 4-shot & 5-shot & 8-shot & 10-shot \\
\midrule
\multirow{2}{*}{Llama-3.1-8B} & GMRID\_v2 & 71.51\% & \textbf{87.08\%} & - & 86.52\% & - & 86.35\% & - & 77.28\% \\
 & GMRID\_v3 & 74.61\% & \textbf{89.50\%} & 81.04\% & - & 80.52\% & - & 79.09\% & 79.20\% \\
\hline
\multirow{2}{*}{Llama-3.1-70B} & GMRID\_v2 & 92.54\% & 93.20\% & - & 93.83\% & - & 93.95\% & - & \textbf{94.01\%} \\
 & GMRID\_v3 & 94.05\% & 95.02\% & 95.25\% & - & 95.16\% & - & 95.23\% & \textbf{95.51\%} \\
\hline
\multirow{2}{*}{gpt-4o-mini} & GMRID\_v2 & 92.18\% & \textbf{93.86\%} & - & 91.92\% & - & 92.54\% & - & 92.97\% \\
 & GMRID\_v3 & 93.68\% & 94.83\% & 94.00\% & - & 94.11\% & - & 94.59\% & \textbf{94.92\%} \\
\hline
\multirow{2}{*}{gpt-4o} & GMRID\_v2 & 94.97\% & 97.42\% & - & 97.46\% & - & \textbf{98.76\%} & - & 97.12\% \\
 & GMRID\_v3 & 95.39\% & 98.27\% & 98.61\% & - & \textbf{99.74\%} & - & 99.00\% & 98.13\% \\
\bottomrule
\end{tabular}
\end{adjustbox}
\caption{F1 Score Comparison Across Few-Shot Settings for \textit{GMRID v2} and \textit{GMRID v3} Datasets}
\label{tab:few_shot_results}
\end{table*}

\subsection{Performance Trade-offs Between Quantized and FP16 Precision LLMs}
\label{sec:quantization}

Ollama models default to \texttt{Q4\_K\_M} quantization, a technique that significantly reduces memory usage and computational overhead. This raises a critical question: how do these quantized models perform compared to their FP16 precision counterparts? To investigate, we evaluated both 4-bit and FP16 variants of models with up to 14 billion parameters on two hardware platforms—the RTX 4090 and M4 Max.

\begin{figure*}[htbp]
    \centering
    \includegraphics[width=.75\textwidth]{images/4bit_vs_fp16.png}
    \caption{Performance Comparison Between 4-bit Quantized Models and FP16 Models Across Hardware Platforms. 
    Subplot (a) displays the performance of various models on the RTX 4090 platform, while subplot (b) presents the corresponding performance on the M4 Max platform. 
    % Blue bars represent the F1 score (percentage), measured on the left y-axis, and red bars denote throughput (tokens per second), measured on the right y-axis. 
    % The x-axis labels each model with its name and the number of shots (in parentheses) that achieved the highest F1 score. 
    % Results for the 4-bit quantized variant and the FP16 variant of each model are shown side by side to facilitate direct comparison, highlighting trade-offs between accuracy and computational efficiency.
    }
    \label{fig:4bit_vs_fp16}
\end{figure*}

Fig.~\ref{fig:4bit_vs_fp16} illustrates the comparative analysis of performance and throughput across these platforms. Specifically, we examined models ranging in size from 0.5 billion to 14 billion parameters, capturing metrics such as F1 score and throughput. 

Key findings include:
\begin{itemize}
    \item Models in 4-bit precision consistently demonstrate faster evaluation times compared to FP16 across both platforms, making them highly efficient for inference tasks.
    
    \item FP16 models occasionally achieve marginally higher F1 scores—such as with smaller models like \textit{qwen2.5:1.5b}, which showed a 7.34\% improvement on the RTX 4090. However, for most models, these differences are often negligible or even negative.
    
    \item 4-bit quantization delivers significant throughput gains on both platforms, reinforcing its suitability for resource-constrained edge AI deployments.
    
    \item For larger models, such as \textit{llama3.2-vision:11b} and \textit{qwen2.5:14b}, FP16 exhibited substantial throughput reductions on both platforms, highlighting the computational efficiency of 4-bit quantization for high-demand tasks.
\end{itemize}


\begin{table}[htbp]
\centering
\begin{adjustbox}{width=0.45\textwidth}
\begin{tabular}{llrrr}
\toprule
\textbf{Model} & \textbf{\makecell{Memory\\(Load/Eval)}} & \textbf{\makecell{\% Mem.\\Reduction\\(Load/Eval)}} & \textbf{F1 (\%)} & \textbf{T-put (t/s)} \\
\midrule
\texttt{qwen2.5:0.5b} & 1.4GB/2.4GB & 30.0\%/20.0\% & 48.20 & 13,123 \\
\texttt{qwen2.5:0.5b (FP16)} & 2.0GB/3.0GB & --- & 46.70 & 12,428 \\
\cmidrule(lr){1-5}
\texttt{llama3.2:1b} & 2.8GB/5.3GB & 30.0\%/18.5\% & 63.50 & 5,035 \\
\texttt{llama3.2:1b (FP16)} & 4.0GB/6.5GB & --- & 63.50 & 4,084 \\
\cmidrule(lr){1-5}
\texttt{qwen2.5:1.5b} & 2.1GB/3.4GB & 51.2\%/39.3\% & 73.70 & 7,419 \\
\texttt{qwen2.5:1.5b (FP16)} & 4.3GB/5.6GB & --- & 81.00 & 5,611 \\
\cmidrule(lr){1-5}
\texttt{llama3.2:3b} & 4.0GB/8.2GB & 53.5\%/31.7\% & 83.60 & 9,766 \\
\texttt{llama3.2:3b (FP16)} & 8.6GB/12GB & --- & 82.20 & 5,654 \\
\cmidrule(lr){1-5}
\texttt{qwen2.5:3b} & 3.1GB/4.9GB & 58.7\%/47.3\% & 91.30 & 8,653 \\
\texttt{qwen2.5:3b (FP16)} & 7.5GB/9.3GB & --- & 90.30 & 6,584 \\
\cmidrule(lr){1-5}
\texttt{qwen2.5:7b} & 6.0GB/9.0GB & 62.5\%/52.6\% & 92.60 & 6,607 \\
\texttt{qwen2.5:7b (FP16)} & 16GB/19GB & --- & 93.00 & 2,491 \\
\cmidrule(lr){1-5}
\texttt{llama3.1:8b} & 6.7GB/11GB & 60.6\%/50.0\% & 92.00 & 6,761 \\
\texttt{llama3.1:8b (FP16)} & 17GB/22GB & --- & 94.10 & 3,805 \\
\cmidrule(lr){1-5}
\texttt{llama3.2-vision:11b} & 12GB/13GB & 52.0\%/50.0\% & 92.50 & 6,293 \\
\texttt{llama3.2-vision:11b (FP16)} & 25GB/26GB & --- & 94.40 & 468 \\
\cmidrule(lr){1-5}
\texttt{qwen2.5:14b} & 11GB/18GB & 64.5\%/52.6\% & 94.90 & 3,995 \\
\texttt{qwen2.5:14b (FP16)} & 31GB/38GB & --- & 94.20 & 89 \\
\cmidrule(lr){1-5}
\texttt{qwen2.5:32b} & 23GB/31GB & N/A & 96.90 & 2,128 \\
\cmidrule(lr){1-5}
\texttt{llama3.1:70b} & 44GB/55GB & N/A & 95.50 & 94 \\
\cmidrule(lr){1-5}
\texttt{llama3.3:70b} & 46GB/58GB & N/A & 95.70 & 73 \\
\cmidrule(lr){1-5}
\texttt{qwen2.5:72b} & 51GB/63GB & N/A & 96.90 & 109 \\
\cmidrule(lr){1-5}
\texttt{llama3.2-vision:90b} & 60GB/63GB & N/A & 96.10 & 56 \\
\bottomrule
\end{tabular}
\end{adjustbox}
\caption{Performance and resource requirements of LLMs. \texttt{FP16} refers to the FP16 precision variant of each model. “\% Mem. Reduction (Load/Eval)” shows memory savings for loading and evaluation with Q4\_K\_M quantization compared to FP16. F1 scores are percentages, and throughput (T-put) is measured in tokens per second (t/s). Metrics were collected on the RTX 4090 platform.}
\label{tab:model_memory_footprints}
\end{table}

An important consideration in the trade-offs between 4-bit quantized and FP16 models lies in their memory efficiency. Table~\ref{tab:model_memory_footprints} illustrates the significant differences in memory usage between these two precision types:

\begin{itemize}
    \item 4-bit quantized models consistently use 30.0--64.5\% less memory during loading compared to FP16. For instance, \textit{qwen2.5:1.5b} requires only 2.1GB in its quantized form versus 4.3GB in FP16. This disparity grows with model size: \textit{qwen2.5:14b} needs 11GB in 4-bit but 31GB in FP16.

    \item During evaluation, 4-bit models reduce memory usage by 20.0--52.6\%. For example, the \textit{qwen2.5:3b} model uses 4.9GB in 4-bit versus 9.3GB in FP16—a nearly 90\% increase. Interestingly, this is higher than the 9.0GB required for the evaluation of the larger \textit{qwen2.5:7b} model in 4-bit. Moreover, despite its higher memory demand, the FP16 \textit{qwen2.5:3b} performs worse, achieving a lower F1 score (90.30\% vs. 92.60\%) and reduced throughput (6,584 vs. 6,607 tokens/s).

    \item The reduced memory footprint of 4-bit models facilitates deployment on hardware with limited VRAM, such as consumer GPUs or edge devices. For instance, the \textit{qwen2.5:7b} model requires 16GB/19GB (Load/Eval) in FP16 but only 6GB/9GB in its quantized form. This makes 4-bit models practical for devices with modest memory capacities while preserving competitive performance metrics.
\end{itemize}

\subsection{Performance Analysis Across All Quantized Models}
\label{sec:all_quantized_models}

Due to hardware limitations, only the RTX A6000 [Windows], RTX A6000, RTX 4090, and M3 Max systems successfully completed the evaluation of all quantized models within a reasonable timeframe. Fig.~\ref{fig:f1_scores_all_models} presents the performance patterns across model sizes on these platforms. As the RTX A6000 [Windows] and RTX A6000 configurations produce nearly identical F1 scores across all models and shot configurations, the RTX A6000 [Windows] is excluded from this analysis.

\begin{figure*}[htbp]
    \centering
    \includegraphics[width=0.8\textwidth]{images/f1_scores_all_models.png}
    \caption{F1 scores of quantized models on RTX A6000, RTX 4090, and M3 Max across different model sizes and shot counts.}
    \label{fig:f1_scores_all_models}
\end{figure*}

Our analysis reveals several key findings:

\paragraph{Small Models (0.5b--3b)} These models exhibit limited scalability with the number of shots:
\begin{itemize}
    \item Models smaller than 1.5b, such as \texttt{qwen2.5:0.5b} and \texttt{llama3.2:1b}, show modest F1 score improvements or even slight decreases across 0--10 shots, highlighting their limited few-shot learning capabilities.
    \item \texttt{qwen2.5:1.5b} achieves a notable F1 score of 73.66--75.13\% at 1 shot but exhibits diminishing returns with additional shots.
    \item Performance patterns remain consistent across all three GPU platforms.
    \item \texttt{qwen2.5:3b} achieves the highest performance among small models, reaching an F1 score of 91.00--91.30\% at 8 shots.
\end{itemize}

\paragraph{Medium Models (7b--14b)} These configurations demonstrate robust scaling:
\begin{itemize}
    \item Significant F1 score improvements are observed with few-shot learning, particularly in the 0--1 shot range.
    \item Performance remains consistent across the %RTX A6000, RTX 4090, and M3 Max 
    platforms, underscoring strong quantization robustness.
    \item \texttt{qwen2.5:32b} leads this category, achieving a remarkable F1 score of 96.68--96.95\% at 8 or 10 shots.
\end{itemize}

\paragraph{Large Models (70b--90b)} These models showcase optimal performance:
\begin{itemize}
    \item Exceptional F1 scores exceeding 95\% are achieved with as few as 4 shots.
    \item Minimal variation across GPU platforms indicates effective quantization scaling for larger models.
    \item Peak performance is achieved by \texttt{qwen2.5:72b}, reaching an impressive F1 score of 96.59--97.08\% at 8 shots.
\end{itemize}

% Across small, medium, and large models, the \texttt{qwen2.5} series consistently achieves the highest scores across these two platforms, signifying the robustness of its architecture and the effectiveness of its quantization strategy.

This analysis reveals that quantization effectively preserves model performance while enabling deployment across diverse GPU architectures. Larger models consistently outperform smaller variants, albeit at the cost of increased computational requirements. The results also highlight the effectiveness of few-shot learning in quantized settings, with most models demonstrating significant improvements within the first four shots.

\subsection{Cross-Platform Performance Analysis and Hardware-Specific Insights}
\label{sec:cross-platform}

Tables~\ref{tab:small-models}, \ref{tab:medium-models}, and \ref{tab:large-models} present performance metrics under few-shot settings that yielded the highest F1 scores across hardware platforms, grouped by model size. Each table reports F1 scores, throughput (tokens/second), and mean evaluation time (seconds), along with their statistical distributions. Due to memory constraints, large models (70B--90B) are evaluated only on high-capacity systems (RTX A6000, RTX A6000 [Windows], RTX 4090, and M3 Max), while all other models are tested on every available platform.

The analysis reveals four key findings:

\begin{enumerate}
    \item Ollama's \texttt{Q4\_K\_M} quantization consistently outperforms bitsandbytes's 4-bit implementation used in baseline experiments. For example, \texttt{llama3.1:8b} achieves a mean F1 score of 92.78\% ($\pm$0.50\%) across six platforms, surpassing its baseline of 89.50\%. Similarly, \texttt{llama3.1:70b} achieves 95.83\% ($\pm$0.07\%) across three platforms, exceeding its baseline of 95.51\% (Table~\ref{tab:few_shot_results}). These results highlight the efficiency and effectiveness of Ollama's quantization strategy.

    \item A performance comparison between WSL and Windows on identical GPUs reveals notable trends. Both configurations yield nearly identical F1 scores across all model sizes, but WSL typically achieves higher throughput and lower evaluation times. For example, on the RTX A6000, \texttt{qwen2.5:0.5b} achieves 16,933 tokens/s with WSL compared to 816 tokens/s on Windows—a 21x difference. This gap narrows with larger models, reducing to a 15\% difference at 90B. On the RTX 4090 Laptop, throughput differences decrease from 9x at 0.5B to 3x at 14B. 
    The only exception is \texttt{qwen2.5:32b} on the RTX 4090 Laptop, which exhibits slightly slower speeds on WSL (221 tokens/s) compared to Windows (228 tokens/s). This may be influenced by its large memory requirements, as the model requires 23GB of memory for loading, far exceeding the GPU's 16GB VRAM capacity. 
    In summary, while WSL provides superior efficiency for small to medium-sized models, it may not perform as well for workloads involving larger models. Given that most consumer-grade NVIDIA GPUs operate on Windows, leveraging WSL is essential for optimal efficiency. However, users must consider its limitations for specific configurations.
    
    \item Model size predictably impacts computational metrics. For instance, on the RTX 4090, throughput decreases from 13,124 tokens/s (\texttt{qwen2.5:0.5b}) to 56 tokens/s (\texttt{llama3.2-vision:90b}), while mean evaluation time rises from 0.16 seconds to 26.505 seconds. This inverse relationship between model size and throughput is observed across all evaluated platforms, with scaling rates varying based on hardware capabilities.

    \item The \texttt{qwen2.5} series consistently achieves superior F1 scores across all model sizes and platforms, reaffirming insights from Section~\ref{sec:all_quantized_models}. Notably, the \texttt{qwen2.5:32b} model achieves an F1 score of 96.75\% ($\pm$0.11\%) while maintaining practical inference speeds, showcasing its balanced optimization of performance and computational efficiency.
\end{enumerate}

% Small Models (0.5b-3b)
\begin{table}[htbp]
    \centering
    \begin{adjustbox}{width=0.4\textwidth}
    \begin{tabular}{llrrr}
\toprule
\textbf{Model} & \textbf{Platform/Stats} & \textbf{F1 (\%)} & \textbf{T-put (t/s)} & \textbf{Time (s)} \\
\midrule
\multirow[t]{10}{*}{qwen2.5:0.5b} & RTX A6000 & 45.85 & 16933 & 0.124 \\
 & RTX A6000 [Windows] & 45.85 & 816 & 2.573 \\
 & RTX 4090 & 48.22 & 13124 & 0.160 \\
 & M3 Max & 47.45 & 7178 & 0.258 \\
 & M4 Max & 46.08 & 9632 & 0.218 \\
 & RTX 4090 Laptop & 46.58 & 6907 & 0.304 \\
 & RTX 4090 Laptop [Windows] & 46.58 & 776 & 2.707 \\
 & Jetson AGX Orin & 46.65 & 1762 & 1.192 \\
 \cmidrule(lr){2-5}
 & mean & 46.66 & 7141 & 0.942 \\
 & std & 0.77 & 5561 & 1.032 \\
\cline{1-5}
\multirow[t]{10}{*}{llama3.2:1b} & RTX A6000 & 63.14 & 3922 & 0.276 \\
 & RTX A6000 [Windows] & 63.14 & 401 & 2.698 \\
 & RTX 4090 & 63.51 & 5035 & 0.215 \\
 & M3 Max & 63.50 & 1899 & 0.570 \\
 & M4 Max & 63.34 & 2400 & 0.451 \\
 & RTX 4090 Laptop & 62.65 & 1893 & 0.572 \\
 & RTX 4090 Laptop [Windows] & 62.65 & 361 & 2.998 \\
 & Jetson AGX Orin & 63.41 & 352 & 3.074 \\
 \cmidrule(lr){2-5}
 & mean & 63.17 & 2033 & 1.357 \\
 & std & 0.33 & 1622 & 1.223 \\
\cline{1-5}
\multirow[t]{10}{*}{qwen2.5:1.5b} & RTX A6000 & 74.38 & 8744 & 0.134 \\
 & RTX A6000 [Windows] & 74.38 & 456 & 2.571 \\
 & RTX 4090 & 73.66 & 7419 & 0.158 \\
 & M3 Max & 75.13 & 2224 & 0.527 \\
 & M4 Max & 74.09 & 4055 & 0.289 \\
 & RTX 4090 Laptop & 74.94 & 3768 & 0.311 \\
 & RTX 4090 Laptop [Windows] & 74.94 & 429 & 2.733 \\
 & Jetson AGX Orin & 72.16 & 740 & 1.583 \\
 \cmidrule(lr){2-5}
 & mean & 74.21 & 3479 & 1.038 \\
 & std & 0.90 & 2983 & 1.028 \\
\cline{1-5}
\multirow[t]{10}{*}{llama3.2:3b} & RTX A6000 & 83.87 & 12220 & 0.171 \\
 & RTX A6000 [Windows] & 83.87 & 803 & 2.601 \\
 & RTX 4090 & 83.61 & 9766 & 0.214 \\
 & M3 Max & 83.40 & 2521 & 0.829 \\
 & M4 Max & 84.33 & 3779 & 0.553 \\
 & RTX 4090 Laptop & 84.18 & 4850 & 0.380 \\
 & RTX 4090 Laptop [Windows] & 84.18 & 662 & 2.783 \\
 & Jetson AGX Orin & 84.34 & 596 & 3.091 \\
 \cmidrule(lr){2-5}
 & mean & 83.97 & 4400 & 1.328 \\
 & std & 0.32 & 4115 & 1.182 \\
\cline{1-5}
\multirow[t]{10}{*}{qwen2.5:3b} & RTX A6000 & 91.00 & 10118 & 0.183 \\
 & RTX A6000 [Windows] & 91.00 & 710 & 2.608 \\
 & RTX 4090 & 91.30 & 8653 & 0.214 \\
 & M3 Max & 91.07 & 1197 & 1.547 \\
 & M4 Max & 91.03 & 3802 & 0.487 \\
 & RTX 4090 Laptop & 90.65 & 4797 & 0.386 \\
 & RTX 4090 Laptop [Windows] & 90.65 & 659 & 2.808 \\
 & Jetson AGX Orin & 90.46 & 730 & 2.536 \\
 \cmidrule(lr){2-5}
 & mean & 90.89 & 3833 & 1.346 \\
 & std & 0.26 & 3537 & 1.089 \\
% \cline{1-5}
\bottomrule
\end{tabular}
    \end{adjustbox}
    \caption{Performance Metrics for Small Models (0.5b-3b). F1 scores are expressed as percentages, throughput (T-put) is measured in tokens per second (t/s), and mean evaluation time (Time) is measured in seconds (s).}
    \label{tab:small-models}
\end{table}

% Medium Models (7b-32b)
\begin{table}[htbp]
    \centering
    \begin{adjustbox}{width=0.4\textwidth}
    \begin{tabular}{llrrr}
\toprule
\textbf{Model} & \textbf{Platform/Stats} & \textbf{F1 (\%)} & \textbf{T-put (t/s)} & \textbf{Time (s)} \\
\midrule
\multirow[t]{10}{*}{qwen2.5:7b} & RTX A6000 & 92.74 & 7298 & 0.288 \\
 & RTX A6000 [Windows] & 92.83 & 784 & 2.681 \\
 & RTX 4090 & 92.65 & 6607 & 0.318 \\
 & M3 Max & 92.50 & 1194 & 1.760 \\
 & M4 Max & 92.72 & 2041 & 1.030 \\
 & RTX 4090 Laptop & 92.46 & 3944 & 0.533 \\
 & RTX 4090 Laptop [Windows] & 92.46 & 705 & 2.983 \\
 & Jetson AGX Orin & 92.78 & 410 & 5.130 \\
 \cmidrule(lr){2-5}
 & mean & 92.64 & 2873 & 1.840 \\
 & std & 0.14 & 2582 & 1.579 \\
\cline{1-5}
\multirow[t]{10}{*}{llama3.1:8b} & RTX A6000 & 93.10 & 7855 & 0.266 \\
 & RTX A6000 [Windows] & 93.10 & 778 & 2.684 \\
 & RTX 4090 & 91.96 & 6761 & 0.309 \\
 & M3 Max & 92.21 & 1430 & 1.461 \\
 & M4 Max & 93.10 & 2269 & 0.921 \\
 & RTX 4090 Laptop & 93.19 & 3935 & 0.531 \\
 & RTX 4090 Laptop [Windows] & 93.19 & 706 & 2.959 \\
 & Jetson AGX Orin & 93.10 & 406 & 5.146 \\
 \cmidrule(lr){2-5}
 & mean & 92.87 & 3017 & 1.785 \\
 & std & 0.46 & 2704 & 1.597 \\
\cline{1-5}
\multirow[t]{10}{*}{llama3.2-vision:11b} & RTX A6000 & 92.63 & 7944 & 0.263 \\
 & RTX A6000 [Windows] & 92.63 & 780 & 2.680 \\
 & RTX 4090 & 92.53 & 6293 & 0.332 \\
 & M3 Max & 93.12 & 1232 & 1.696 \\
 & M4 Max & 92.64 & 2261 & 0.924 \\
 & RTX 4090 Laptop & 93.39 & 3927 & 0.532 \\
 & RTX 4090 Laptop [Windows] & 93.39 & 707 & 2.956 \\
 & Jetson AGX Orin & 92.82 & 456 & 4.585 \\
 \cmidrule(lr){2-5}
 & mean & 92.89 & 2950 & 1.746 \\
 & std & 0.33 & 2657 & 1.447 \\
\cline{1-5}
\multirow[t]{10}{*}{qwen2.5:14b} & RTX A6000 & 95.18 & 4105 & 0.512 \\
 & RTX A6000 [Windows] & 95.00 & 728 & 2.888 \\
 & RTX 4090 & 94.88 & 3995 & 0.526 \\
 & M3 Max & 94.78 & 343 & 3.418 \\
 & M4 Max & 94.84 & 1033 & 2.035 \\
 & RTX 4090 Laptop & 95.17 & 1988 & 1.057 \\
 & RTX 4090 Laptop [Windows] & 95.09 & 638 & 3.292 \\
 & Jetson AGX Orin & 95.06 & 204 & 10.304 \\
 & mean & 95.00 & 1629 & 3.004 \\
 & std & 0.14 & 1487 & 2.971 \\
 \cmidrule(lr){2-5}
 & mean & 94.99 & 1591 & 6.191 \\
 & std & 0.16 & 1523 & 9.086 \\
\cline{1-5}
\multirow[t]{10}{*}{qwen2.5:32b} & RTX A6000 & 96.68 & 2048 & 0.906 \\
 & RTX A6000 [Windows] & 96.68 & 560 & 3.315 \\
 & RTX 4090 & 96.86 & 2128 & 0.872 \\
 & M3 Max & 96.95 & 259 & 8.131 \\
 & M4 Max & 96.86 & 143 & 12.995 \\
 & RTX 4090 Laptop & 96.68 & 221 & 8.404 \\
 & RTX 4090 Laptop [Windows] & 96.60 & 228 & 8.130 \\
 & Jetson AGX Orin & 96.68 & 95 & 19.497 \\
 & mean & 96.75 & 710 & 7.781 \\
 & std & 0.11 & 806 & 5.916 \\
% \cline{1-5}
\bottomrule
\end{tabular}
    \end{adjustbox}
    \caption{Performance Metrics for Medium Models (7b-32b). F1 scores are expressed as percentages, throughput (T-put) is measured in tokens per second (t/s), and mean evaluation time (Time) is measured in seconds (s).}
    \label{tab:medium-models}
\end{table}

% Large Models (70b-90b)
\begin{table}[htbp]
    \centering
    \begin{adjustbox}{width=0.4\textwidth}
    \begin{tabular}{llrrr}
\toprule
\textbf{Model} & \textbf{Platform/Stats} & \textbf{F1 (\%)} & \textbf{T-put (t/s)} & \textbf{Time (s)} \\
\midrule
\multirow[t]{6}{*}{llama3.1:70b} & RTX A6000 & 95.96 & 621 & 2.374 \\
 & RTX A6000 [Windows] & 95.96 & 312 & 4.724 \\
 & RTX 4090 & 95.46 & 94 & 15.657 \\
 & M3 Max & 95.55 & 120 & 12.293 \\
 \cmidrule(lr){2-5}
 & mean & 95.73 & 287 & 8.762 \\
 & std & 0.23 & 211 & 5.411 \\
\cline{1-5}
\multirow[t]{6}{*}{llama3.3:70b} & RTX A6000 & 95.87 & 937 & 1.966 \\
 & RTX A6000 [Windows] & 95.87 & 426 & 4.322 \\
 & RTX 4090 & 95.73 & 73 & 20.241 \\
 & M3 Max & 95.90 & 37 & 39.523 \\
 \cmidrule(lr){2-5}
 & mean & 95.84 & 368 & 16.513 \\
 & std & 0.07 & 362 & 15.030 \\
\cline{1-5}
\multirow[t]{6}{*}{qwen2.5:72b} & RTX A6000 & 97.08 & 515 & 3.604 \\
 & RTX A6000 [Windows] & 97.08 & 314 & 5.911 \\
 & RTX 4090 & 96.85 & 109 & 17.072 \\
 & M3 Max & 96.59 & 116 & 16.055 \\
 \cmidrule(lr){2-5}
 & mean & 96.90 & 263 & 10.660 \\
 & std & 0.20 & 167 & 5.970 \\
\cline{1-5}
\multirow[t]{6}{*}{llama3.2-vision:90b} & RTX A6000 & 96.05 & 117 & 12.611 \\
 & RTX A6000 [Windows] & 96.05 & 102 & 14.460 \\
 & RTX 4090 & 96.13 & 56 & 26.505 \\
 & M3 Max & 95.93 & 85 & 17.379 \\
 \cmidrule(lr){2-5}
 & mean & 96.04 & 90 & 17.739 \\
 & std & 0.07 & 23 & 5.339 \\
% \cline{1-5}
\bottomrule
\end{tabular}
    \end{adjustbox}
    \caption{Performance Metrics for Large Models (70b–90b). F1 scores are expressed as percentages, throughput (T-put) is measured in tokens per second (t/s), and mean evaluation time (Time) is measured in seconds (s).}
    \label{tab:large-models}
\end{table}

\subsection{Discussions}

Our analysis, integrating results from Sections~\ref{sec:quantization}, \ref{sec:all_quantized_models}, and \ref{sec:cross-platform}, reveals distinct performance patterns across model sizes, each presenting unique trade-offs for edge deployment.

\paragraph{Small Models (0.5B--3B Parameters)} 
\begin{itemize}
    \item High performance variability (standard deviation: 0.26--0.90\%) underscores significant hardware dependencies.
    \item Models with fewer than 3B parameters exhibit limited few-shot capabilities. For instance, \texttt{qwen2.5:0.5b} achieves a mean F1 score of only 46.66\%; however, performance improves markedly with increased scale (\texttt{qwen2.5:3b}: 90.89\% F1). 
    \item Task-specific fine-tuning appears to be more effective for enhancing performance, though further investigation is necessary to confirm its efficacy.
\end{itemize}

\paragraph{Medium Models (7B--32B Parameters)} 
\begin{itemize}
    \item These models achieve substantial F1 score gains in early shots (0--1), with steady improvements observed in subsequent shots.
    \item Low performance variance (standard deviation: 0.11--0.46\%) across hardware platforms underscores their potential for reliable deployment.
    \item Employing 4-bit quantization reduces memory usage by 50.0--52.6\% without compromising performance metrics.
\end{itemize}

\paragraph{Large Models (70B--90B Parameters)} 
\begin{itemize}
    \item Achieve F1 scores exceeding 95\% within four shots, underscoring efficient few-shot learning capabilities.
    \item Significant memory and computational demands constrain deployment flexibility.
    \item Minimal F1 score variability (standard deviation: 0.07--0.23\%) ensures consistent cross-platform performance.
\end{itemize}

These findings underscore the importance of scale-specific optimization strategies. Small models benefit from platform-sensitive fine-tuning, medium models offer a practical balance for edge deployments, and large models provide superior accuracy at the cost of increased resource consumption.

\section{Conclusion}

This study evaluates LLMs across diverse hardware platforms using the Ollama framework, offering valuable insights into edge AI deployments. Ollama’s \texttt{Q4\_K\_M} quantization consistently outperforms bitsandbytes’ implementation, with \texttt{llama3.1:8b} and \texttt{llama3.1:70b} achieving 3.28\% and 0.32\% higher F1 scores, respectively. The framework’s robust cross-platform stability further underscores its suitability for edge environments.

Our findings reveal a clear relationship between model size and deployment demands. Larger models exhibit improved cross-platform stability, with F1 score variances dropping from 0.26–0.9\% (small models) to 0.07–0.22\% (large models), albeit at higher computational cost. Medium-sized models (7B–32B) strike a balance, exemplified by \texttt{qwen2.5:32B} achieving a 96.75\% ($\pm$0.11\%) F1 score at practical inference speeds.

A comparison between WSL and Windows reveals nearly identical F1 scores, but highlights WSL’s superior efficiency, delivering 9–21× higher throughput for smaller models and over 2× for most larger models. Given that most consumer-grade NVIDIA GPUs run on Windows, we strongly recommend using WSL for optimal performance.

The findings and recommendations from this study enhance the understanding of deploying LLMs at the edge, enabling more efficient, scalable, and accessible AI solutions. By balancing performance, resource utilization, and deployment constraints, this work lays the groundwork for innovative strategies supporting privacy-preserving, low-latency applications. As edge computing evolves, these insights will guide the development of more adaptable and advanced AI systems.

To further advance edge AI, we propose the following research directions:

\begin{enumerate}
    \item Developing scale-aware quantization and fine-tuning for resource-constrained environments.
    \item Optimizing deployments on cost-effective hardware (e.g., NVIDIA’s Jetson Orin Nano Super) and consumer-grade devices.
    \item Establishing standardized benchmarks and real-world deployment case studies for robust evaluation.
\end{enumerate}

These efforts will foster scalable, efficient, and accessible edge AI deployments, enabling privacy-preserving, low-latency applications.


\bibliographystyle{IEEEtran}
\bibliography{references}

\end{document}
